% !TeX program = xelatex
% UTF-8 transcription of the 12 source photographs; not an updated medical guideline.
% Source order (original filenames, preserved JPEGs):
% pic01.jpg = 20261013_065236766.jpg
% pic02.jpg = 20261013_065243617.jpg
% pic03.jpg = 20261013_065249632.jpg
% pic04.jpg = 20261013_065255432.jpg
% pic05.jpg = 20261013_065301459.jpg
% pic06.jpg = 20261013_065308868.jpg
% pic07.jpg = 20261013_065314044.jpg
% pic08.jpg = 20261013_065322432.jpg
% pic09.jpg = 20261013_065328304.jpg
% pic10.jpg = 20261013_065335223.jpg
% pic11.jpg = 20261013_065341158.jpg
% pic12.jpg = 20261013_065346049.jpg
% Compile directly, or input this file into an existing XePersian document.
\newif\ifKDisordersStandalone
\ifdefined\settextfont
  \KDisordersStandalonefalse
\else
  \KDisordersStandalonetrue
\fi
\ifKDisordersStandalone
  \documentclass[11pt,a4paper]{article}
  \usepackage[margin=1.8cm]{geometry}
  \usepackage{amsmath,amssymb,enumitem}
  \usepackage{xepersian}
  \IfFontExistsTF{Amiri}{\settextfont{Amiri}}{\settextfont{FreeSerif}}
  \setlatintextfont{FreeSerif}
  \begin{document}
\fi
\begingroup
\emergencystretch=3em
\ifdefined\setlist
  \setlist[itemize]{nosep,topsep=4pt}
  \setlist[enumerate]{nosep,topsep=4pt}
\fi
\newcommand{\KDen}[1]{\lr{#1}}
\newcommand{\KDpage}[2]{\clearpage\section*{صفحهٔ #1 ـ #2}}

\section*{صفحهٔ ۱ ـ فیزیولوژی و اختلالات پتاسیم}
\begin{itemize}
\item غلظت پتاسیم خارج سلولی: \KDen{$4.2\pm0.3\,\mathrm{mEq/L}$}.
\item ۹۸\% پتاسیم بدن داخل سلول و ۲\% آن بیرون سلول است.
\item حجم مایع داخل سلولی (\KDen{ICF}) معادل ۴۰\% \KDen{TBW}؛ حجم مایع خارج سلولی (\KDen{ECF}) معادل ۲۰\% \KDen{TBW}.
\item در یک فرد ۷۰ کیلوگرمی، \KDen{$3920\,\mathrm{mEq}$} پتاسیم داخل سلول و \KDen{$59\,\mathrm{mEq}$} پتاسیم در مایع خارج سلولی وجود دارد. در حالی که هر وعدهٔ غذایی در حدود \KDen{$50\,\mathrm{mEq}$} پتاسیم دارد.
\item مقدار پتاسیم دفع‌شده از مدفوع فقط ۵ تا ۱۰\% پتاسیم خورده‌شده است؛ عامل عمدهٔ دفع پتاسیم، کلیه است.
\item بیشتر پتاسیم خورده‌شده به سرعت وارد سلول‌ها می‌شود تا بعد به تدریج از کلیه‌ها دفع شود.
\end{itemize}
\subsection*{عوامل شیفت پتاسیم}
عواملی که سبب شیفت پتاسیم به داخل سلول می‌شوند:
\begin{itemize}
\item انسولین؛ آلدوسترون؛ تحریک \KDen{$\beta_2$} آدرنرژیک؛ آلکالوز.
\end{itemize}
عواملی که سبب شیفت پتاسیم به خارج سلول می‌شوند:
\begin{itemize}
\item کمبود انسولین (\KDen{DM})؛ کمبود آلدوسترون (آدیسون).
\item بلوک \KDen{$\beta$} آدرنرژیک؛ اسیدوز؛ ورزش شدید؛ لیز سلولی.
\item افزایش اسمولالیتهٔ مایع خارج سلولی.
\end{itemize}
\begin{itemize}
\item انسولین تزریقی نیز سبب شیفت پتاسیم به داخل سلول می‌شود.
\item افزایش پتاسیم (هایپرکالمی) موجب ترشح آلدوسترون و سپس شیفت پتاسیم به داخل سلول می‌شود.
\end{itemize}

\KDpage{۲}{تنظیم پتاسیم و دفع کلیوی}
\begin{itemize}
\item سندرم کان: همراه هایپوکالمی رخ می‌دهد. \KDen{Conn's syndrome}: ترشح بیش از حد آلدوسترون.
\item کاتکول‌آمین‌ها (خصوصاً اپی‌نفرین)، به دنبال ترشح، سبب شیفت داخل سلولی پتاسیم می‌شوند؛ از طریق فعال‌شدن گیرنده‌های \KDen{$\beta_2$} آدرنرژیک.
\item پروپرانولول: هایپرکالمی.
\item در اسیدوز متابولیک: افزایش غلظت یون هیدروژن (\KDen{$H^+$})، کاهش فعالیت پمپ \KDen{Na--K--ATPase}، کاهش خروج سدیم از سلول و کاهش ورود پتاسیم به داخل سلول؛ در نتیجه هایپرکالمی.
\item طی ورزش شدید، هایپرکالمی به دلیل خروج پتاسیم از عضلات ایجاد می‌شود؛ در افراد تمرین‌نکرده، با بتابلوکر و با کمبود انسولین، بیشتر می‌شود.
\item در موارد نادر می‌تواند منجر به آریتمی قلبی گردد.
\end{itemize}
دفع کلیوی پتاسیم توسط ۳ عامل تعیین می‌گردد:
\begin{enumerate}
\item نرخ فیلتراسیون پتاسیم (\KDen{GFR} \(\times\) غلظت پلاسمایی پتاسیم).
\item میزان بازجذب پتاسیم به وسیلهٔ توبول‌ها.
\item میزان ترشح پتاسیم توسط توبول‌ها.
\end{enumerate}
\begin{equation*}
4.2\,\mathrm{mEq/L}\times180\,\mathrm{L/day}=756\,\mathrm{mEq/day}
\end{equation*}
\begin{itemize}
\item کاهش شدید \KDen{GFR} می‌تواند منجر به هایپرکالمی شدید گردد.
\item در توبول پروگزیمال (\KDen{PCT}): بازجذب ۶۵\% پتاسیم فیلترشده؛ \KDen{$491\,\mathrm{mEq/day}$}.
\item لوپ هنله، خصوصاً قسمت ضخیم صعودکننده (\KDen{TAL}): بازجذب ۲۷\% (۲۵ تا ۳۰\%) پتاسیم فیلترشده؛ \KDen{$204\,\mathrm{mEq/day}$}.
\item بازجذب همراه با سدیم و کلر توسط \KDen{Na--K--2Cl}.
\end{itemize}

\KDpage{۳}{توبول دیستال و سلول اصلی}
\begin{itemize}
\item مقدار پتاسیم بازجذب‌شده در دو مورد قبلی معمولاً در طول روز ثابت است.
\item مقدار بازجذب‌شده در توبول‌های جمع‌کننده وابسته به دریافت روزانهٔ پتاسیم متغیر است: ۴\% (\KDen{$30\,\mathrm{mEq/day}$}).
\item تغییرات عمدهٔ تنظیمی بدن عمدتاً از طریق تغییر در ترشح پتاسیم در \KDen{DCT} و توبول جمع‌کننده (\KDen{CT}) رخ می‌دهد: ۸\% (\KDen{$60\,\mathrm{mEq/day}$}). مهم‌ترین سلول‌ها: سلول‌های اصلی \KDen{DCT} و \KDen{CT}.
\item در افرادی که رژیم پرپتاسیم دارند، دفع پتاسیم از میزان فیلترشده بیشتر است.
\item کلیه‌ها می‌توانند پتاسیم دفع‌شده را تا ۱\% پتاسیم فیلترشده کاهش دهند (تا کمتر از \KDen{$10\,\mathrm{mEq}$} در روز).
\item اگر دریافت روزانهٔ پتاسیم کمتر از \KDen{$10\,\mathrm{mEq}$} باشد، هایپوکالمی ظرف چند روز رخ می‌دهد.
\end{itemize}
\subsection*{سلول اصلی (\KDen{Principal cell})}
سلول اصلی طی دو مرحله پتاسیم را به داخل توبول ترشح می‌کند:
\begin{enumerate}
\item پمپ \KDen{Na--K--ATPase} در غشای بازولترال سلول اصلی، \KDen{K} را وارد سلول و \KDen{Na} را از سلول خارج می‌کند.
\item انتشار غیرفعال پتاسیم به داخل توبول به دلیل غلظت بالای آن در سلول، به وسیلهٔ دو کانال در غشای اپیکال:
  \begin{itemize}
  \item \KDen{ROMK}: \KDen{Renal Outer Medullary K}.
  \item \KDen{BK}: \KDen{high conductance Big K}.
  \end{itemize}
\end{enumerate}
با دریافت زیاد پتاسیم، تعداد این دو کانال زیاد می‌شود.
\subsection*{عوامل کنترل‌کنندهٔ ترشح پتاسیم}
\begin{itemize}
\item فعالیت پمپ \KDen{Na--K--ATPase}.
\item گرادیان الکتروشیمیایی از خون به توبول.
\item نفوذپذیری غشای اپیکال سلول اصلی.
\end{itemize}

\KDpage{۴}{سلول‌های اینترکاله و آلدوسترون}
\begin{itemize}
\item سلول اینترکاله هم می‌تواند پتاسیم را ترشح و هم بازجذب کند؛ دو نوع دارد: \KDen{aIC} و \KDen{bIC}.
\item \KDen{aIC}: بازجذب پتاسیم در شرایط هیپوکالمی؛ در شرایط عادی نقش خاصی ندارد.
\item ترانسپورتر \KDen{H--K--ATPase}: پتاسیم داخل و هیدروژن خارج؛ سپس پتاسیم با انتشار به داخل خون می‌رود.
\item \KDen{bIC}: ترشح پتاسیم در شرایط پتاسیم اضافی؛ پمپ \KDen{H--K} به جای غشای لومینال، در غشای بازولترال قرار دارد.
\end{itemize}
نقش اصلی تنظیم پتاسیم بر عهدهٔ سلول \KDen{principal} است. مهم‌ترین عوامل مؤثر بر این سلول:
\begin{enumerate}
\item افزایش غلظت پتاسیم خارج سلولی: تحریک سلول اصلی و ترشح پتاسیم.
\item افزایش آلدوسترون.
\item افزایش سرعت جریان توبولی.
\end{enumerate}
\subsection*{افزایش غلظت پتاسیم خارج سلولی}
\begin{itemize}
\item تحریک پمپ \KDen{Na--K--ATPase} و تسهیل انتقال \KDen{K} به داخل سلول اصلی.
\item افزایش گرادیان پتاسیم از مایع میان‌بافتی کلیه به داخل سلول اصلی توبول؛ کاهش بک‌لیک پتاسیم از داخل سلول اصلی.
\item افزایش سنتز کانال‌های \KDen{BK} و \KDen{ROMK}.
\item تحریک فوق کلیه، منجر به ترشح آلدوسترون.
\end{itemize}
\subsection*{آلدوسترون}
\begin{itemize}
\item افزایش ترشح پتاسیم از دو طریق:
  \begin{itemize}
  \item افزایش بازجذب فعال \KDen{Na} در سلول‌های اصلی به وسیلهٔ \KDen{ENaC} در غشای اپیکال.
  \item پمپ \KDen{Na--K--ATPase} در غشای بازولترال، \KDen{Na} را خارج و \KDen{K} را وارد سلول اصلی می‌کند.
  \end{itemize}
\item ترشح از طریق کانال‌های \KDen{ROMK} و \KDen{BK}؛ افزایش تعداد کانال‌ها.
\end{itemize}

\KDpage{۵}{جریان توبولی و توزیع پتاسیم}
\begin{itemize}
\item آلدوسترون علاوه بر اثر تحریکی بر ترشح کلیوی پتاسیم، برداشت سلولی پتاسیم را نیز افزایش می‌دهد.
\item افزایش سرعت جریان در توبول دیستال، ترشح پتاسیم را تحریک می‌کند. مثال: افزایش حجم، دریافت مقدار زیاد سدیم، دیورتیک‌ها.
\item تعداد کانال‌های \KDen{BK} و \KDen{ROMK} را افزایش می‌دهد.
\item اسیدوز حاد، ترشح پتاسیم را کاهش می‌دهد: کاهش فعالیت پمپ \KDen{Na--K--ATPase}.
\item در اسیدوز مزمن (چند روزه)، ترشح پتاسیم افزایش می‌یابد: مهار بازجذب \KDen{NaCl} و آب در \KDen{PCT} و افزایش جریان در توبول دیستال.
\item بیشترین میزان پتاسیم در بافت عضله است.
\item بیشترین میزان \KDen{K} خارج سلولی در استخوان است؛ ۱\% کل \KDen{K} بدن در پلاسما قرار دارد.
\item انسولین فعالیت پمپ \KDen{Na--K--ATPase} را تسهیل می‌کند؛ اثر \KDen{$\beta$} آدرنرژیک نیز همین است.
\item هایپرکالمی، ترشح انسولین را تحریک می‌کند.
\item آگونیست \KDen{$\alpha$} آدرنرژیک و ورزش: خارج‌شدن پتاسیم از سلول.
\item با هر \KDen{$10\,\mathrm{mOsm}$} افزایش اسمولالیتهٔ پلاسما، پتاسیم سرم \KDen{$0.6\,\mathrm{mEq/L}$} افزایش می‌یابد.
\item دوز توصیه‌شدهٔ دریافت پتاسیم در رژیم غذایی، \KDen{$1\text{--}2\,\mathrm{mEq/kg}$} در روز است.
\item حدود ۹۰\% پتاسیم خورده‌شده جذب می‌شود؛ اکثر جذب در پروگزیمال رودهٔ کوچک است.
\item نارسایی کلیه، آلدوسترون و گلوکوکورتیکوئیدها، ترشح \KDen{K} در کولون را افزایش می‌دهند.
\item بیش از ۴۰\% پتاسیم هر وعدهٔ غذایی (\KDen{acute K load})، بلافاصله بر اثر انسولین و اپی‌نفرین وارد سلول می‌شود.
\end{itemize}

\KDpage{۶}{هایپرکالمی و هایپرکالمی کاذب}
\begin{itemize}
\item گلوکوکورتیکوئیدها و \KDen{ADH}: فلو دیستال افزایش می‌یابد، \KDen{$Na^+$} دیستال افزایش می‌یابد؛ افزایش دفع ادراری پتاسیم.
\item انسولین، کاتکول‌آمین‌ها و آنژیوتانسین \KDen{II}: کاهش دفع ادراری پتاسیم.
\item \KDen{ADH}: افزایش دفع ادراری پتاسیم و افزایش بازجذب آب؛ در مجموع تأثیر خاصی بر پتاسیم ندارد.
\end{itemize}
\subsection*{هایپرکالمی}
اتیولوژی هایپرکالمی: افزایش دریافت \KDen{K}، شیفت به بیرون سلول و کاهش دفع.
\subsection*{سودوهایپرکالمی}
\begin{itemize}
\item همولیز در نمونهٔ خون رخ می‌دهد. در کودکان شایع است؛ تورنیکه بسته نشود.
\item سطح پتاسیم سرم تقریباً \KDen{$0.4$} بالاتر از سطح پتاسیم پلاسما است، به دلیل آزادشدن \KDen{K} طی تشکیل لخته.
\item به ازای افزایش هر \KDen{$100{,}000$} پلاکت، پتاسیم سرم \KDen{$0.15$} بالا می‌رود.
\item اندازه‌گیری سطح پتاسیم پلاسما بهتر و دقیق‌تر است.
\item نگهداری نمونه در دمای سرد: آزادشدن \KDen{K} از سلول‌ها، سودوهایپرکالمی.
\item در دمای اتاق: جذب \KDen{K} توسط سلول‌ها، سودوهایپوکالمی.
\item فامیلیال سودوهایپرکالمی (\KDen{AD}): نفوذ \KDen{K} از غشای \KDen{RBC}ها نشت پیدا می‌کند.
\end{itemize}
\subsection*{شیفت به خارج سلول}
\begin{itemize}
\item اسیدوز تنفسی اثر خیلی کمتری بر شیفت خارج سلولی \KDen{K} نسبت به اسیدوز متابولیک دارد.
\item اسیدهای ارگانیک (لاکتیک اسید، کتواسید) اثر ضعیف‌تری دارند.
\item در بیهوشی بیمار با فلج هایپوکالمی نباید از سوکسینیل‌کولین استفاده کنیم؛ شیفت \KDen{K} به خارج سلول ایجاد می‌کند.
\end{itemize}

\KDpage{۷}{علل هایپرکالمی و اختلالات ارثی}
\begin{itemize}
\item فلج دوره‌ای هایپرکالمیک (\KDen{AD}): نقص کانال \KDen{$Na^+$} و رهاسازی دوره‌ای پتاسیم از سلول.
\item نوزادان، خصوصاً نوزادان پره‌ترم، به دلیل عدم بلوغ کامل کلیه، در خطر افزایش‌یافتهٔ هایپرکالمی هستند.
\item هر اختلال آدرنال (ارثی، اکتسابی) که منجر به کاهش تولید آلدوسترون شود، منجر به هایپرکالمی می‌گردد.
\item بیماران به طور تیپیک اسیدوز متابولیک، سالت‌ویستینگ و هایپوناترمی نیز دارند.
\item فرم شایع هایپرپلازی مادرزادی آدرنال (\KDen{CAH}): کمبود ۲۱ هیدروکسیلاز.
\item شیرخوار پسر: هایپرکالمی، اسیدوز متابولیک، هایپوناترمی و علائم دهیدریشن.
\item دختران با این اختلال به دلیل ژنیتالیای مبهم در بدو تولد تشخیص داده می‌شوند.
\item هایپورنینمی در بیماری‌های کلیه: تولید آلدوسترون کم؛ در این بیماران اسیدوز متابولیک و هایپرکالمی بدون هایپوناترمی.
\item اختلال‌های مربوط به توبول کلیه (\KDen{RTAs}): اختلال ترشح پتاسیم به ادرار و هایپرکالمی.
\end{itemize}
\subsection*{سودوهایپوآلدوسترونیسم تیپ \KDen{I}}
\begin{itemize}
\item هایپرکالمی، اسیدوز متابولیک و سالت‌ویستینگ (کلرید سدیم) شدید.
\item فرم \KDen{AR}: نقص کانال \KDen{$Na^+$}؛ هایپوناترمی، علائم دهیدریشن و آلدوسترون بالا.
\item علائم: \KDen{FTT}، اسهال، تهوع، راش پوستی؛ شروع در شیرخوارگی.
\item فرم \KDen{AD}: نقص گیرندهٔ آلدوسترون؛ خفیف‌تر، شروع در بزرگسالی.
\end{itemize}
\subsection*{سودوهایپوآلدوسترونیسم تیپ \KDen{II}}
\begin{itemize}
\item شامل هایپرکالمیک هایپرتانسیون (سندرم \KDen{Gordon})، \KDen{AD}.
\item \KDen{HTN} به دلیل ریتنشن نمک و اختلال دفع \KDen{K}؛ وابسته به هایپرکالمی.
\item اسیدوز هایپرکلرمیک متابولیک.
\item درمان: تیازید دیورتیک.
\end{itemize}

\KDpage{۸}{داروها، تظاهرات بالینی و ارزیابی هایپرکالمی}
\begin{itemize}
\item سندرم بارتر: نقص \KDen{ROMK} (بارتر تیپ ۲)؛ هایپرکالمی گذرای نوزادی، در ادامه هایپوکالمی.
\item \KDen{ACEI}: تشدید هایپرکالمی در بیمار آنوریک با مهار دفع گوارشی پتاسیم.
\item دیورتیک‌های پتاسیم‌اسپرینگ به سادگی منجر به هایپرکالمی می‌شوند، چون این بیماران اغلب در حال دریافت مکمل خوراکی پتاسیم هستند.
\item \KDen{OCP}های حاوی \KDen{Drospirenone} اثر آلدوسترون را مهار می‌کنند؛ در بیمار با اختلال عملکرد کلیه، به سادگی هایپرکالمی ایجاد می‌کنند.
\end{itemize}
\subsection*{تظاهرات بالینی}
\begin{itemize}
\item شروع تغییرات \KDen{EKG}: \KDen{peaked T wave}، \KDen{STD}، افزایش \KDen{PR interval}، \KDen{flat P wave} و \KDen{wide QRS}.
\item ارتباط بین سطح پتاسیم و یافتهٔ \KDen{EKG} ضعیف است.
\item در نهایت \KDen{VF}.
\item پارستزی، ضعف عضلانی و فلج بالارونده؛ بعد از تغییرات قلبی.
\end{itemize}
\subsection*{آزمایش‌های درخواستی}
\begin{itemize}
\item \KDen{VBG}، \KDen{Cr}، \KDen{BUN}.
\item افزایش \KDen{CPK} همراه هایپوکلسیمی: \KDen{K} از سلول مرده آزاد می‌شود؛ رابدومیولیز، برنز و \KDen{TLS}.
\item هایپرفسفاتمی و هایپراوریسمی همزمان؛ در لکوز نیز دیده می‌شود.
\end{itemize}
\subsection*{درمان}
\begin{itemize}
\item اخذ \KDen{EKG} اگر \KDen{$K>6.5$}.
\item قطع تمام پتاسیم‌های \KDen{IV} و \KDen{PO}.
\item در صورت نیاز به \KDen{PC}، از \KDen{Washed RBCs} استفاده شود.
\end{itemize}

\KDpage{۹}{درمان هایپرکالمی و علل هایپوکالمی}
اجزای اصلی درمان:
\begin{enumerate}
\item پایدارسازی قلب و جلوگیری از آریتمی.
\item خارج‌کردن \KDen{K} از بدن.
\end{enumerate}
\begin{itemize}
\item کلسیم غشا میوسیت‌ها را پایدار می‌کند؛ در حالت عادی ظرف چند دقیقه تزریق می‌شود (\KDen{IV}).
\item در بیمار تحت درمان با دیژیتال، ظرف \KDen{$30\,\mathrm{min}$}؛ کلسیم آریتمی ایجاد می‌کند.
\item بیکربنات شیفت پتاسیم به داخل سلول ایجاد می‌کند؛ اگر بیمار به طور همزمان اسیدوز متابولیک داشته باشد، گزینهٔ مناسبی است.
\item انسولین + گلوکز: ظرف ۳۰ دقیقه پتاسیم را به داخل سلول شیفت می‌دهد.
\item نبولایزر آلبوترول: تحریک گیرندهٔ \KDen{$\beta_1$} آدرنرژیک و شیفت داخل سلول. % beta_1 is as written in the source.
\item اگر بیمار آنوریک نیست، دیورتیک لوپ می‌دهیم.
\item \KDen{SPS} (سدیم پلی‌استیرن سولفونات)، \KDen{Kayexalate}: خوراکی یا رکتال.
\item دیالیز اگر نارسایی شدید یا تولید بی‌امان هایپرکالمی مثل \KDen{TLS} وجود دارد.
\end{itemize}
\subsection*{هایپوکالمی}
چهار مکانیزم اصلی دارد:
\begin{itemize}
\item کاذب: در بیمار لوکمی با \KDen{WBC} خیلی بالا، اگر نمونه مانده باشد در دمای اتاق.
\item شیفت به داخل سلول.
\item کاهش \KDen{intake}.
\item \KDen{loss} رنال یا اکسترارنال.
\end{itemize}

\KDpage{۱۰}{علل و تظاهرات بالینی هایپوکالمی}
\begin{itemize}
\item سندرم سودوبارتر در بیمار \KDen{CF} در هوای گرم، به دلیل تعریق: هایپوکالمیک هایپوکلرمیا آلکالوز.
\item آلکالمی یک علت اصلی شیفت پتاسیم به داخل سلول است؛ این اثر در آلکالوز متابولیک بیشتر از تنفسی است.
\item مسمومیت با تولوئن (استنشاق رنگ و چسب): هایپوکالمی بر اثر شیفت.
\item فلج دوره‌ای هایپوکالمیک (\KDen{AD}): سلول‌ها به صورت حاد پتاسیم را برداشت می‌کنند.
\item ورزش شدید در هوای گرم: مقدار قابل توجه \KDen{$K^+$} از طریق عرق از دست می‌رود.
\item \KDen{Urinary K wasting} ممکن است همراه با اسیدوز متابولیک باشد (\KDen{RTA} پروگزیمال یا دیستال).
\end{itemize}
\subsection*{تظاهرات بالینی هایپوکالمی}
\begin{itemize}
\item \KDen{EKG}: \KDen{flattened T}، \KDen{STD}، \KDen{U wave}.
\item \KDen{VF}، \KDen{TdP}: معمولاً بیماری قلبی زمینه‌ای وجود دارد.
\item هایپوکالمی قلب را به شدت نسبت به آریتمی در زمینهٔ دیژیتال حساس می‌کند؛ \KDen{SVT}، \KDen{VT} و بلوک‌ها.
\item ضعف و کرامپ عضلانی نیز رخ می‌دهد.
\item فلج در \KDen{$K<2.5$} رخ می‌دهد؛ از پاها شروع و به سمت بالا می‌رود.
\item در صورت فلج تنفسی: \KDen{IMV}.
\item کاهش حرکات \KDen{GI}: یبوست؛ در \KDen{$K<2.5$} ایلئوس رخ می‌دهد.
\item اختلال عملکرد مثانه: احتباس ادرار.
\item پلی‌اوری و پلی‌دیپسی: \KDen{DI} نفروژنیک.
\item تحریک تولید آمونیاک در کلیه؛ خصوصاً اگر نارسایی کبد وجود داشته باشد، هایپرآمونمی می‌تواند ایجاد شود.
\item تشدید انسفالوپاتی کبدی.
\end{itemize}

\KDpage{۱۱}{تشخیص و اصول درمان هایپوکالمی}
\begin{itemize}
\item هایپوکالمی مزمن به کلیه آسیب می‌زند: کیست کلیوی و نفریت بینابینی.
\end{itemize}
\subsection*{تشخیص}
\begin{itemize}
\item بررسی رژیم غذایی، \KDen{GI loss} و داروها.
\item وجود \KDen{HTN}: وجود اثر زیاد مینرالوکورتیکوئیدها.
\item وجود اسیدوز متابولیک: اسهال، \KDen{RTA} دیستال و پروگزیمال.
\item وجود آلکالوز متابولیک: استفراغ، \KDen{NGT}، دیورتیک، آلدوسترون بالا و سندرم‌های بارتر و جیتلمن.
\item اگر اتیولوژی پیدا نشد: اندازه‌گیری \KDen{K} ادراری جهت افتراق \KDen{loss} کلیوی از غیرکلیوی.
\item جمع‌آوری ادرار ۲۴ ساعته؛ نسبت \KDen{K} به \KDen{Cr} در ادرار \KDen{spot}؛ نسبت ترشح پتاسیم.
\item محاسبهٔ ترانس‌توبولار پتاسیم گرادیان (\KDen{TTKG}).
\end{itemize}
\begin{equation*}
\mathrm{TTKG}=\frac{U_K}{P_K}\times\frac{P_{\mathrm{osmol}}}{U_{\mathrm{osmol}}}
\end{equation*}
\begin{itemize}
\item حتماً باید اسمولالیتهٔ ادرار بیشتر از سرم باشد تا بتوان از \KDen{TTKG} استفاده کرد.
\item \KDen{$\mathrm{TTKG}>4$} در حضور هایپوکالمی، بیانگر \KDen{loss} پتاسیم کلیوی است.
\end{itemize}
\subsection*{درمان}
\begin{itemize}
\item موارد شدید یا علامت‌دار نیازمند درمان تهاجمی هستند.
\item در موارد نارسایی کلیه، مکمل پتاسیم با دقت تجویز شود، چون کلیه توان دفع پتاسیم اضافی را ندارد.
\item وقتی ۱) اسیدوز متابولیک و ۲) کمبود انسولین در \KDen{DKA} وجود دارد، سطح پلاسمایی پتاسیم بیانگر دفیسیت واقعی بیمار نیست. این موارد که اصلاح شوند، دفیسیت واقعی معلوم می‌گردد.
\item در اختلالاتی که فقط شیفت ایجاد شده، پتاسیم را در زمان فلج دوره‌ای با دوز خیلی دقیق جایگزین می‌کنند تا بیمار دچار هایپرکالمی نشود. % This sentence continues at the bottom of pic10.
\end{itemize}

\KDpage{۱۲}{جایگزینی پتاسیم و نکات تکمیلی}
\begin{itemize}
\item فلج دوره‌ای تیروتوکسیک: به سرعت به درمان با پروپرانولول جواب می‌دهد و نیاز به پتاسیم ندارد.
\item پتاسیم خوراکی بسیار ایمن‌تر از وریدی است.
\item دوز خوراکی معمول: \KDen{$1\text{--}2\,\mathrm{mEq/kg/day}$}؛ ماکسیمم \KDen{$60\,\mathrm{mEq/day}$} در دوزهای منقسم.
\item دوز وریدی معمولی: \KDen{$0.5\text{--}1\,\mathrm{mEq/kg}$} که ظرف \KDen{$2\,\mathrm{h}$} تزریق می‌شود؛ حداکثر \KDen{$40\,\mathrm{mEq}$}.
\item پتاسیم کلراید محلول انتخابی است.
\item اسیدوز + هایپوکالمی: پتاسیم سیترات، پتاسیم استات.
\item هایپوفسفاتمی + هایپوکالمی: پتاسیم فسفات.
\item هایپوکالمی + آلکالوز متابولیک + علائم دهیدریشن (\KDen{gastric loss}): جبران حجم با \KDen{N/S}، دفع کلیوی پتاسیم را کاهش می‌دهد.
\item اصلاح هایپومنیزیمی همزمان ضروری است.
\end{itemize}
\subsection*{نکات تکمیلی}
\begin{itemize}
\item اسهال معمولاً هایپوکالمی + اسیدوز متابولیک ایجاد می‌کند، به دلیل دفع بیکربنات در مدفوع.
\item در مصرف مسهل‌ها: آلکالوز متابولیک خفیف یا اسید ـ باز نرمال.
\item \KDen{RTA} پروگزیمال و دیستال: \KDen{Urinary K wasting} + اسیدوز متابولیک.
\item در فاز پلی‌اوریک \KDen{ATN} و دیورز \KDen{postobstructive}، به هر دو صورت رنال \KDen{K wasting} رخ می‌دهد؛ ممکن است با اسیدوز متابولیک همراه باشد.
\item کمبود ایزولهٔ منیزیوم: \KDen{renal K wasting}.
\item پنی‌سیلین یک آنیون است و دفع ادراری \KDen{K} را زیاد می‌کند؛ تنها با نمک سدیم پنی‌سیلین رخ می‌دهد، نه با نمک پتاسیم.
\end{itemize}

\endgroup
\ifKDisordersStandalone
  \expandafter\enddocument
\fi
